%----------------------------------------------------------------------------------------
%	PACKAGES AND OTHER DOCUMENT CONFIGURATIONS
%----------------------------------------------------------------------------------------
\documentclass{report}
\usepackage[utf8]{inputenc}
\usepackage{titlesec}
\usepackage{lmodern}
\usepackage[T1]{fontenc}
\usepackage{fancyhdr} % Required for custom headers
\usepackage{lastpage} % Required to determine the last page for the footer
\usepackage{extramarks} % Required for headers and footers
\usepackage[usenames,dvipsnames]{color} % Required for custom colors
\usepackage{graphicx} % Required to insert images
\usepackage{listings} % Required for insertion of code
\usepackage{courier} % Required for the courier font
\usepackage{lipsum} % Used for inserting dummy 'Lorem ipsum' text into the template
\usepackage[section]{placeins} %when placing img
\usepackage{url}

% Margins
\topmargin=-.45in
\evensidemargin=0in
\oddsidemargin=0in
\textwidth=6.5in
\textheight=9.0in
\headsep=0.25in
\linespread{1.1} % Line spacing

% Set up the header and footer
\pagestyle{fancy}
\lhead{\hmwkAuthorName} % Top left header
\chead{} % Top center head
\rhead{\hmwkTitle} % Top right header
\lfoot{\hmwkClass} % Bottom left footer
\cfoot{} % Bottom center footer
\rfoot{Page\ \thepage\ of\ \protect\pageref{LastPage}} % Bottom right footer
\renewcommand\headrulewidth{0.4pt} % Size of the header rule
\renewcommand\footrulewidth{0.4pt} % Size of the footer rule
\setlength\parindent{0pt} % Removes all indentation from paragraphs

%----------------------------------------------------------------------------------------
%	CODE INCLUSION CONFIGURATION
%----------------------------------------------------------------------------------------



%----------------------------------------------------------------------------------------
%	DOCUMENT STRUCTURE COMMANDS
%	Skip this unless you know what you're doing
%----------------------------------------------------------------------------------------
\definecolor{RoyalRed}{RGB}{157,16, 45}
\titleformat{\chapter}[display]
  {\normalsize \huge \color{black}}
  {\flushright \normalsize \color{RoyalRed} \MakeUppercase { \hspace{1 ex}}  { \fontsize{70}{70}\selectfont \color{RoyalRed} \thechapter}} {-29 pt}{\bfseries\huge}
  \titlespacing*{\chapter}{0pt}{-35pt}{20pt}
  

% you can add \chaptertitlename before hspace

%----------------------------------------------------------------------------------------
%	NAME AND CLASS SECTION
%----------------------------------------------------------------------------------------

\newcommand{\hmwkTitle}{SON - Self-Organizing Networks } % Assignment title
\newcommand{\hmwkDueDate}{Saturday,\ December\ 9,\ 2017} % Due date
\newcommand{\hmwkClass}{Introduction to mobile communication -\ 34330 \\} % Course/class
\newcommand{\hmwkClassTime}{Department of Photonics Engineering} % Kursusansvarlig profession
\newcommand{\hmwkClassInstructor}{Henrik Lehrmann Christiansen,} % Teacher/lecturer
\newcommand{\hmwkAuthorName}{Frederik Rander Andersen, s164146} % Your name

%----------------------------------------------------------------------------------------
%	TITLE PAGE
%----------------------------------------------------------------------------------------

\title{
\vspace{-0.5in}
\noindent\makebox[\linewidth]{\rule{\textwidth}{1pt}} 
\textmd{\textbf{\hmwkClass \hmwkTitle}}\\
\normalsize\vspace{0.1in}\small{\hmwkDueDate}\\
\vspace{0.1in}\large{\textit{\hmwkClassInstructor\ \hmwkClassTime}}
\noindent\makebox[\linewidth]{\rule{\textwidth}{1pt}} 
\vspace{1in}
}

\author{\textbf{\hmwkAuthorName}}
\date{} % Insert date here if you want it to appear below your name

%----------------------------------------------------------------------------------------

\begin{document}
\begin{figure}
	\centering
	\includegraphics[scale=0.15]{dtulogo.png}
\end{figure}

\begin{figure}[!b]
	\includegraphics[scale=0.25]{dtufooter.png}
\end{figure}
\maketitle

%----------------------------------------------------------------------------------------
%	TABLE OF CONTENTS
%----------------------------------------------------------------------------------------

\newpage
\tableofcontents
\newpage

%----------------------------------------------------------------------------------------
%	INTRODUCTION & OVERVIEW
%----------------------------------------------------------------------------------------

\chapter{SON - Introduction and overview}
\section{Introduction}
This report will be taking a closer look at SON (Self-Organizing Networks) and what it means for the industry of telecommunications. 
The features of SON aims to improve end user experience and reduce the costs entailed with providing a network, while still increasing the quality and efficiency of the network. 
These features, along with their impact will be explored later in the report 

\section{Overview}
All mobile networks need to be managed and as systems become more and more complex, the need for better and easier ways to manage them are important as ever. LTE (Long Term Evolution) is the newest technology and also the most complex. Therefore, in LTE, management needs to be as good as possible. SON (Self-Organizing Networks) is a very promising area for providers, as it makes network-management cheaper, more efficient and easier. This is also why SON is most prevalent in LTE networks, simply because the demands of LTE are much higher and therefore LTE networks are quite complex.\\ The goal of SON is basically to reduce the need for technicians and increase the network capabilities, such that the network will be as good as possible in regards to coverage, capacity and user experience. Generally, SON has three main areas; self-configuration, self-optimization and self-healing. These will discussed in depth later. 

\section{Why SON?}
The reasons for using SON are very obvious from a provider standpoint. First of all, the cost of a Self-Organizing Network should be much lower than a traditional network because you simply don't need as many technicians and other workers associated with configuration and maintenance of a network. Furthermore, the network operator will be able to better monitor, improve and repair the network, since the entities in the network automatically measure performance and also adjust based on what is needed. From a end-user standpoint SON is also a good thing, since it will mean better coverage and capacity, while keeping cost down since the network operator is not spending as much money on the network. Simply put, SON will provide a better network, cheaper and more efficiently than before.
	%----------------------------------------------------------------------------------------
%	SELF-CONFIGURATION
%----------------------------------------------------------------------------------------
\chapter{Self-configuration}
\section{Main idea and overview}
The first area of SON is self-configuration. The main idea behind the self-configuration part of SON is to automate the setup of eNBs (eNodeB). This allows a plug and play type of setup, which saves the network owner a lot of time and money, since you would usually need a technician to setup new eNBs, which could take a lot of time.  
The self-configuration also reduces the risk of incorrect installation and integration of eNBs into the existing network. The amount of needed cells is also rising with the increase in network usage.

\section{Features of self-configuration}
There are three main features of self-configuration in LTE; self-configuration of eNB, Automatic Neighbour Relations and automatic configuration of Physical Cell ID (PCI). 
Another small but important feature of self-configuration is Dynamic Radio Configuration (DNC). 
In order for the eNB to configure itself correctly, it will need to alter the planned data a little. This is done automatically by the eNB itself by performing measurements and thus adjusting initial power, antenna tilt etc. %citer http://www.radio-electronics.com/info/cellulartelecomms/self-organising-networks-son/self-configuration.php

\subsection{Process of eNB self-configuration}

One of the key ideas of self-configuration is self-configuration of a new eNB trying to connect to the network. The eNB is in this case not connected to anything but the network management subsystem and Serving-Gateway (SGW). The following steps are performed in order to connect the new eNB to the network. 

\begin{itemize}
	\item First, the eNB is powered on and plugged in where necessary, then it will have an already 	established connection, which it will use until the radio frequency transmission is turned on.
	\item The DNS/DHCP server will now provide an IP address to the eNB.
	\item Now, the self-configuration subsystem of the of operation and management information is 		sent to the eNB. 
	\item A gateway is now configured to the eNB such that it can now connect and communicate with other internet nodes through IP packets.
	\item The eNB will provide its own information e.g. hardware, ID, supported technologies etc. to the self-configuration subsystem in order to get identified and authenticated. 
	\item The eNB will now be able to download the correct software and radio-configuration information. 
	\item After this download, the eNB configures itself according to the downloaded transport and radio configuration information. 
	\item Now, the eNB can connect to the Operation Administration Management (OAM) for other management functions.
	\item The S1 interface is then established, giving the eNB connection to the Evolved Packet Core Network (EPCN). The X2 interface is also established and the eNB is now connected to other eNBs in the network. \cite{Atayero2014}
\end{itemize}

\subsection{Automatic Neighbour Relations (ANR)}

Automatic Neighbour Relation (ANR), is the process of managing the neighbour relations, such that the cells know who their neighbouring cells are and what technologies they support. It is crucial that this information is updated and is correct, otherwise handovers might fail and thereby result in dropped calls. If the neighbouring cells support different technologies e.g. one might support HSPA and its neighbour LTE, the user needs to know this so it know what frequencies to listen on. \\ Generally the ANR process looks a bit like this:

\begin{itemize}
	\item UE detects unknown PCI and reports it to the serving eNB by sending a Radio Resource Controller (RRC) reconfiguration message.
	\item The serving eNB will now request the UE to send E-UTRAN Cell Global ID (ECGI) of the unknown eNB. The UE then does this by reading the BCCH. 
	\item The serving eNB will now retrieve the IP address from the MME based on the ECGI of the unknown eNB. This allows the serving eNB to correctly setup the X2 interface, since that hasn't been done, because the unknown eNB is unknown. 
	\item The functions will now be extended to cases of inter-Radio Access Technology (inter-RAT) and inter-frequency. \cite{Atayero2014}
\end{itemize}

A detailed figure of a possible ANR procedure can be seen on figure 2.1.

\begin{figure}[!ht]
	\centering
	\includegraphics[scale=0.6]{/home/frederik/Pictures/anr1.png}
	\caption{ANR process \cite{Roessler}}
\end{figure}

On the figure, we see the whole process of ANR, with the UE first detecting the unknown eNB, then sending measurements to the serving eNB, which will then ask for the ECGI of the unknown eNB. The UE will retrieve that information and relay it to the serving eNB which then updates its neighbour cell list and sends that list to the OAM. 
\subsection{Automatic configuration of Physical Cell ID (PCI)}
Every cell in the network will have its own PCI, which is in the SCH, to be used when synchronizing with UE on downlink. In the E-UTRAN there are only 504 PCIs, therefore they have to be reused, however they need to be unique within specific regions. There are two rules which has to be obeyed when configuring PCIs: 
Neighbouring cells must not share PCI. 
All neighbours of one cell must have different PCIs. \cite{Feng2008}


%----------------------------------------------------------------------------------------
%	SELF-OPTIMIZATION
%----------------------------------------------------------------------------------------
\chapter{Self-optimization}
\section{Main idea and overview}
After the network has self-configured it should also be able to adapt according to load and other factors. This is where self-optimization comes in. To ensure effective management and operation of the network, the network should, after configuring itself, have algorithms and processes to make decisions in regards to efficiency and performance. If implemented correctly this will increase performance and stability while keeping cost and energy at a minimum. This is done by taking measurements from both UEs and eNBs and using these to regulate the network. 

\section{Features of self-optimization}
There are several important features of self-optimization, which all contribute to making SON much better. 

\subsection{Mobility Load Balancing (MLB)}
Mobile Load Balancing (MLB) is especially relevant in cases where a cell with high load and a cell with low load are neighbours. Here MLB uses automated algorithms and functions to avoid cell overloading and consequent degradation of performance. This is achieved by having the algorithms adjust parameters to balance the load between cells, these algorithms also prevent other issues that could arise in these situations e.g. ping pong handover. \\
Simply put, MLB in eNBs tries to balance load relatively evenly between neighbouring cells, improve network capacity by regulating cell congestions and manage the network efficiently such that performance is as high as possible. Figure 3.1 shows the basic principle of MLB. 
\begin{figure}[!ht]
	\centering
	\includegraphics[scale=0.3]{/home/frederik/Pictures/hej1234.png}
	\caption{MLB basic principle \cite{Hahn2014}}
\end{figure}
\\
MLB consists of three main parts when used between cells; load reporting, load balancing action and amending mobility configuration.
Here is an example of how MLB could work:
\begin{itemize}
	\item eNB A detects that it is being overloaded
	\item eNB A sends a Resource Status Request to its neighbour eNB B
	\item eNB B answers with a Resource Status response
	\item eNB B sends a Resource Status Update
	\item eNB A will now find suitable UEs to handover to eNB B
	\item Now the handover procedure begins
	\item eNB A sends a Mobility Change Request to eNB B (to change eNB As range)
	\item eNB B answers with a Mobility Change ACK
	\item both eNBs now updates their handover settings to ensure that the UEs won't go back to the previously overloaded eNB A \cite{Sartori2012}
\end{itemize}

\subsection{Mobility Robustness Optimization (MRO)}
Mobility Robustness Optimization (MRO) is mainly used to guarantee good mobility for UE and thereby giving the end-users a good experience. Furthermore, it is used to ensure that handovers are as seamless as possible. Previously, a network operator would have to manually adjust handover parameters in cells, which is very costly, MRO automates this process, meaning lower costs for the network operator and better experience for the end-user. \\ It is important to mention that call re-establishment is only possible in LTE, so if the target cell only supports different Radio Access Technologies (RATs), it will not be possible to restore the connection. 
Another thing that MRO tries to do is minimizing handover ping pong. While ping pong does not lead to RLF, it is however, an inefficient use of resources and could reduce the users throughput. \cite{Sartori2012}

\subsubsection{Backward Handover}
Backward handover occurs when handover information is exchanged between the eNB and the UE on the old radio path, instead of the new radio path between the target eNB and the UE. Then the old eNB, also called source eNB, prepares the target eNB for handover. So, in backward handover, it is the network that perform cell switching and notifies the UE of the target eNB. 
Backward handover is used when the Radio Frequency conditions are declining. \cite{Qualcomm2010}

\subsubsection{Radio Link Failure (RLF) handover}
The way MRO achieves all this is by minimizing call drops and RLF. Sometimes a connection can be restored after a RLF but before the call is dropped. This is known as RLF handover, and occurs when a backward handover partially fails. The process of RLF handover looks a bit like this: 
\begin{itemize}
	\item UE detects radio link problems and starts its RLF timer.
	\item RLF timer expires and the UE now searches for a new target eNB and it now attempts to re-establish the connection with the new eNB. This is all done while the UE is in a connected state. 
	\item The connection is restored if the UE managed to send a measurement report to the old eNB, such that the target eNB could be prepared.\cite{Qualcomm2010}  %citer qualcomm whitepaper
\end{itemize}
It should be noted that no data is lost in this procedure due to the eNB forwarding all data contained in its buffer. 
Sometimes, the mobility problems are so severe that the call is dropped. 

\subsection{Minimization of Drive Tests (MDT)}
Testing is very important for network operators, as it helps to collect statistics about the network. A lot of the testing done on the the network is drive tests, which consists of sending out engineers in a car, then they drive around and test the network in different areas. This requires a lot of special tools, it also takes a long time and is quite costly. Therefore, MDT is very important to network operators. The idea behind MDT is to use the UEs to measure and log information along with location, just like the drive tests. Because of the amount of users and their different behaviour this method of testing would be very thorough, and would give the network operators an exact view of what their customers are experiencing. The main use of MDT would be to optimize network coverage e.g. by detecting network "holes". \cite{Sartori2012}

\subsection{Energy Savings}
Another promising area of self-optimization is energy savings. By reducing power consumption, a network operator can save money and be less taxing on the environment. In SON networks, energy consumption can be reduced both in the UE and inside the network, however, this report will focus on the energy savings that can be made within the network.

\subsubsection{Dynamic cell activation}
In SON, it is possible to dynamically activate and deactivate cells, based on their load, such that you can save power by having low-loaded cells turned off. 
The process can be implemented by letting the MME manage the activation of these cells. This works by having the UE inform the MME, when it is moving close to a cell. \\ The MME would now turn on the cell and when the UE gets within range, it will connect to the cell. When the UE then moves out of the cell, and there is not any other devices connected, the UE will send a message to the MME and the MME will then be able to turn off the cell, thus saving power. \\ Of course, when doing this, one would have to expand coverage of other cells such that coverage is not reduced. \cite{Feng2008}

\subsubsection{Transmission power management}
Based on UE measurements, the transmission power of an eNB can be adjusted, such that it is never using more power than necessary. This needs to be coordinated with load balancing and the general coverage of the network, such that there is not coverage "holes" and cells don't become overloaded because another cell is transmitting with less power. \cite{Feng2008}

\subsection{Coverage and Capacity Optimization (CCO)}
Coverage and Capacity Optimization is used to make sure that the network is performing as intended i.e. making sure that the coverage and capacity is on par with the targets set by the operator. The CCO algorithm finds the optimal antenna configuration and radio frequency settings.
The dilemma CCO faces is coverage vs. capacity, since it is not possible to optimize for both, it has to carefully balance the two, such that the network is working as optimal as possible. If it only focused on capacity, users would experience coverage "holes" and if it focused only on coverage, the capacity might be too slow for a satisfying user experience. 
CCO gets better as time goes on since it constantly gathers data on the network. \cite{Atayero2014}

\subsection{Inter-Cell Interference Coordination (ICIC)}
In LTE networks interference between cells is not uncommon. This leads to poor signal quality, if not dealt with properly. One way of handling it is by using an Inter-Cell Interference Coordination algorithm. \\ This kind of algorithm aims to reduce inter-cell interference so that  radio signal quality will be as high as possible. This is done by giving all cells in area information about the other cells in that same area, so they can calculate what frequencies to use. This coordination happens in each of the cells. The process will be a bit like this: 
First, a cell will collect information about other cells and store that in its Neighbour Relation Table (NRT), then it will assign itself a frequency, which it will attempt to use, the frequency should not already be in use nearby. \\ Now, the cells will use ICIC to optimize their small local area, in which interference could be a problem. The algorithm will calculate the optimal frequency for each cell in the local area and assign them accordingly. \\ When the aforementioned happens, it can at times lead to ping-pong behaviour, where the cells will switch between two or more frequencies. However, the algorithm will then detect this behaviour and choose an arrangement that will break this "loop". \cite{Karla2010} \\ A small area network using ICIC could something like the one pictured on figure 3.2, where the colours represent different frequencies. In the first case, we see interference where the two cells coverage meet, but when using ICIC (on the right) we see that there are two different frequencies in the edge cases, which minimizes interference.

\begin{figure}[!ht]
	\centering
	\includegraphics[scale=0.4]{/home/frederik/Pictures/hejhej.png}
	\caption{Example of ICIC in LTE \cite{Henrik}}
\end{figure}

\subsection{RACH optimization}
The Random Access Channel (RACH) is used for uplink synchronization between UE and the network, initial access, handover etc. So, RACH is very important as it is used to synchronize the UEs and the network. 
The performance of RACH is therefore critical to having a well-working network. The performance is dependent on several factors such as number of UEs in cell, number of handovers, number of calls in cell etc. 
Depending on how the network is configured, meaning transmission power, load, antenna tilt etc., these factors will be affected.\cite{Feng2008} The aforementioned configuration parameters are controlled by other automated processes in self-optimization and self-configuration, so the RACH will see performance increases if these other SON features are working well. \\ In summation, it can be said that RACH optimization will provide reduced connection time, higher user throughput, better cell coverage and network capacity. \\ In order to do this, all measurements from the UE and eNB are sent to the SON components in the eNB. The eNB can now provide this information over the X2 interface to other eNBs in the area. This way the eNBs can use information from neighbours and themselves in order to do RACH optimization. \cite{Atayero2014}


%----------------------------------------------------------------------------------------
%	SELF-HEALING
%----------------------------------------------------------------------------------------
\chapter{Self-healing}
\section{Main idea and overview}
Network maintenance is not only about optimizing the network and adapting it to different scenarios, it is also about making sure that the network is not experiencing severe downtime. Downtime, loss of service or lack of availability are all very dangerous issues that can occur in networks, and it might lead to customers leaving the network, if these problems occur too frequently. This can lead to a network operator losing money because of failures in the network. Self healing therefore uses automatic detection and localization of failures in order to figure out what solution could restore the lost functionality. 

\section{Features of self-healing}
There are some smaller features that self-healing consists of, but the main features are cell outage detection and compensation and cell degradation detection and management. 

\subsection{Cell Outage Detection and Compensation}
A cell outage happens when an eNB becomes unavailable to the network, because of hardware or software issues. It can also happen during network maintenance or if some power saving plan turns off the eNB. There are ways to compensate for this, e.g. by having nearby cells increase transmit power and "expanding" their coverage to the area of the now unavailable eNB. The adjacent eNBs might adjust their antenna tilt in order to be able to cover the area of the unavailable eNB. \\ When an eNB fails, it is critical to recover and re-establish the network as quickly as possible, so that the users don't become dissatisfied. \\ It is not always possible to recover from loss of coverage by increasing antenna tilt and transmit power, as one could imagine a scenario where these parameters are already maximized for coverage e.g. out in the country, where the need for capacity is very low. \\ In cases, where it is possible to recover from loss of coverage however, the coverage and capacity conditions will not be optimal. This should only be temporary, as the network operator should repair the unavailable eNB e.g. by sending out a technician to fix the unavailable eNB. \cite{Sartori2012}

\subsection{Cell degradation detection}
The term cell degradation is used when the difference in expected cell performance and actual cell performance is so significant that the cell is no longer considered healthy. There are several "levels" of degradation ranging from relatively small performance drops to complete outages. The "fix" for each of these levels might differ, as an outage can require on-site repair, whereas some performance issues can be fixed by a remote cell reset. \\
There are several functions that an operator might use to detect cell degradation. 
\subsubsection{Counters}
One of the methods to detect cell degradation is using counters in eNB software to keep track of issues, this could be a fault counter that counts when a major fault happens inn the cell. If these counters get too high, the network operator might consider solving the underlying issues.
\subsubsection{Key Performance Indicators (KPI)}
KPI utilizes the counters in order to calculate the relative performance of a cell. Key Quality Indicators (KQI) uses KPIs to give an broader overview of network performance. It is from these parameters that a network operator decides to fix issues. 
There are several performance indicators which are used to calculate KPI and KQI. These might come from many different sources including: eNBs, UEs, neighbour eNBs etc. 
\subsubsection{Alarms}
Alarms are set to trigger if certain events happen e.g. cell outage, power loss etc. This way the operator quickly knows if something is wrong. Often the alarms are triggered by critical events that might result in severe network performance.  \cite{snov}

%----------------------------------------------------------------------------------------
%	LTE releases
%----------------------------------------------------------------------------------------
\chapter{SON features in LTE releases}
SON features have not always been in LTE and they haven't been added at the same time. Some SON features may not even be implemented yet, but here is an short overview of when some of the SON features discussed in this report was introduced in LTE. 
\subsection{LTE release 8}
LTE release 8 (2006) introduced eNB automatic software download and configuration of X2 and S1. It also included ANR and some support for MLB.
\subsection{LTE release 9}
Release 9 (2008) was an improvement to release 8, and therefore many of the things included in release 9 was optimizations to features of release 8. Features like load balancing optimization, MRO, RACH optimization, CCO and ICIC were part of release 9. 
\subsection{LTE release 10}
Release 10 included further enhancements to the features already present from release 8 and 9. Cell Outage Detection, enhanced ICIC (eICIC) and minimization of drive testing were main parts of release 10. 
\subsection{LTE release 11}
LTE release 11 included even further enhancements to MRO and energy savings. Inter-RAT energy saving management was also introduced.\cite{Atayero2014} \cite{3gpp}

%----------------------------------------------------------------------------------------
%	Conclusion 
%----------------------------------------------------------------------------------------
\chapter{Conclusion}
SON introduces a lot of very nice and convenient features such as eNB self-configuration, ANR, automatic configuration of PCI, MLB, MRO, MDT, energy savings, CCO, ICIC, RACH optimization, Cell outage detection and compensation, Cell degradation detection etc. These features are the most important features of SON and they are constantly getting improved as the technology advances. \\ Self-configuration makes it possible for a network operator to setup eNBs with much less trouble than before and much faster and cheaper. Self-configuration also makes eNBs aware of their neighbours and makes updating neighbour lists very easy with the introduction of ANR. \\ Self-optimization makes the whole network very efficient, while having great performance. This is done by introducing MLB, MRO, MDT, energy savings, CCO, ICIC and RACH optimization. Self-optimization has a lot of features and it is also one of the most complex parts of SON, however, it is also arguably almost the most important, as it brings a lot improvements to the network.  \\ Self-healing is the last feature of SON that this report discussed, and it is also very important as it makes the robustness of the network much higher. It also provides very important information to the network operator, such as network quality, which gives the operator a better understanding of the end-user experience. \\ SON is a very new technology that is starting to grow into the networks more and more, and could very well revolutionary in telecommunication networks. 

%----------------------------------------------------------------------------------------
%	BIBLIOGRAPHY
%----------------------------------------------------------------------------------------
\clearpage
\addcontentsline{toc}{chapter}{References}
\bibliographystyle{ieeetr}
\bibliography{/home/frederik/Documents/bibtex/SON.bib}
\end{document}